\section{Introduction}

Geospatial foundation models (GeoFMs) have recently emerged as a promising paradigm for Earth observation analysis. Models such as Prithvi-100M~\cite{jakubik2023prithvi}, SatMAE~\cite{cong2022satmae}, Scale-MAE~\cite{reed2023scalemae}, and related self-supervised transformers~\cite{hong2024spectralgpt} suggest that large-scale pre-training on satellite imagery can yield transferable representations for downstream remote sensing tasks. At the same time, practitioners increasingly face adaptation problems that differ substantially from the homogeneous settings used during pre-training. Real-world inputs may contain different band subsets, missing channels, alternative sensor combinations, or strongly out-of-distribution modality configurations.

In mainstream natural language processing and computer vision, parameter-efficient fine-tuning (PEFT) methods are widely used to adapt large pre-trained models at low cost. LoRA~\cite{hu2022lora}, adapters~\cite{houlsby2019adapter}, BitFit~\cite{zaken2022bitfit}, and prompt-style tuning methods~\cite{lester2021prompt,jia2022vpt} have become standard choices because they often approach full fine-tuning performance with far fewer trainable parameters. This success naturally raises the question of whether the same PEFT rankings hold for geospatial foundation models.

We argue that this assumption is non-trivial for at least three reasons. First, geospatial foundation models often operate on heterogeneous multi-band inputs rather than fixed RGB inputs. Second, remote sensing pre-training objectives, especially masked autoencoding on spectral-spatial data, differ from the language modeling and image classification setups where PEFT methods are most established. Third, implementation details that are harmless in conventional vision backbones can become critical in GeoFMs, especially when architectural components such as fused attention projections interact with adaptation layers.

Despite the practical importance of this problem, current literature lacks a systematic benchmark of PEFT methods under heterogeneous geospatial input settings. Existing studies often evaluate only one or two adaptation strategies, focus on a single downstream task, or treat PEFT as an implementation detail rather than the object of study. As a result, practitioners still lack clear guidance on which PEFT method to choose when adapting a model such as Prithvi-100M to non-standard remote sensing inputs.

A second, related gap is that PEFT comparisons are almost always reported on curated public datasets. It is rarely demonstrated that the same ranking reproduces on operational GIS pipelines, where imagery is heterogeneous in sensor and resolution, labels are derived from third-party products of varying quality, and validation splits must respect spatial structure to prevent leakage. The absence of such a demonstration is a real obstacle for practitioners deciding whether to ship a PEFT-adapted GeoFM into a production workflow. A third gap, complementary to the first two, is that within-domain PEFT comparisons cannot tell a practitioner what to expect when the target domain drifts substantially from pretraining. Cross-domain protocols on remote-sensing benchmarks --- where train and test are drawn from explicitly disjoint imaging conditions, geographies, or sensor families --- are routine in the dataset literature but largely unstudied for PEFT method ranking on geospatial foundation models.

This paper addresses all three gaps through a controlled empirical study on Prithvi-100M, paired with a production-style validation and a public cross-domain replication. We benchmark seven adaptation strategies across five modality configurations on EuroSAT, validate the main findings on BigEarthNet-S2, extend the protocol to semantic segmentation on LandCover.ai, re-instantiate the resulting five-method comparison on a 35{,}920-patch Linhe County corpus collected from 203 Chinese-domestic high-resolution scenes (0.42--4.29\,m GSD, 22 satellite families) with Esri 2022 global LULC used as supervisory labels, and finally replicate the same five-method protocol on the LoveDA urban/rural cross-domain benchmark. Beyond reporting final scores, we include a full fine-tuning ceiling, a diagnostic split-QKV LoRA ablation, a split-QKV rank-sensitivity audit, a completed EuroSAT channel-bridge ablation, a synthetic weak-label control on the Linhe corpus, and a per-class diagnostic on the LoveDA cross-domain split. The Linhe control is not an independent manual validation set; it tests whether the measured adapter delta changes when the same imagery is paired with a deliberately simple weak label. Across these experiments, we show that PEFT rankings commonly assumed in mainstream NLP and vision should not be transferred to Prithvi-100M without architecture-aware validation. We also show that benchmark rankings can be stress-tested under a geographic-split production-style workflow, while cross-domain segmentation exposes capacity and label-source limits that are not visible in within-domain classification alone.

Our main contributions are as follows:
\begin{enumerate}[leftmargin=1.5em]
    \item We report a systematic benchmark of PEFT methods on Prithvi-100M under heterogeneous input settings, covering 109 public-benchmark experiments across EuroSAT, BigEarthNet-S2, and LandCover.ai.
    \item We separate implementation failure from post-fix capacity limits for LoRA on a fused-QKV Prithvi-100M backbone: standard module-wise insertion can miss packed attention projections, while corrected split-QKV low-rank adaptation still saturates below bottleneck adapters in the tested settings.
    \item We introduce an architecture-aware PEFT audit protocol that records trainable tensors, tests fused-attention insertion, checks channel-bridge sensitivity, and compares low-rank capacity against bottleneck-adapter capacity before interpreting negative PEFT results.
    \item We show that Houlsby adapters are the strongest method in our Prithvi-100M benchmark, achieving 94.5\% of full fine-tuning performance with only 1.4\% of the trainable parameters on EuroSAT and a $+155\%$ relative mIoU gain on dense prediction.
    \item We confirm across datasets that modality selection has a larger practical impact than PEFT selection: richer spectral inputs consistently outperform reduced-band alternatives.
    \item We provide a bounded result for standalone input-stage adaptation (GeoAdapter): it is unstable on public classification benchmarks near Prithvi's pretraining modality, but becomes useful on RGB-only production-style inputs where the modality gap is larger.
    \item We test the benchmark ranking on a Linhe County LULC workflow under a geographic scene-level split, recovering Houlsby $+0.116$ mIoU over linear probing with Esri-derived supervisory labels; we also provide a synthetic weak-label control showing that the measured PEFT delta is sensitive to label source, while noting that this is not an independent manual validation set.
    \item We replicate the five-method ranking on LoveDA urban/rural cross-domain segmentation across 30 runs and report evidence consistent with an adapter-capacity threshold under domain shift. We treat this threshold as a hypothesis supported by the present diagnostics, not as a universal rule for all GeoFMs.
    \item We release a reproducible Geo-MLOps benchmarking framework together with the AlphaEarth-System platform that orchestrates the Linhe deployment, including its data pipeline, training engine, and web dashboard.
\end{enumerate}
